\chapter{Z3 Discharge and Counterexample Search}
\label{ch:z3}

The \Zthree{} \citep{Z3} backend is the second-most-trusted
discharge channel, after \Lean{}~4's kernel.  This chapter explains
which formulas the backend handles, how counterexample search
operates, and what the
\texttt{AXIOM\_TRUSTED}~$\to$~\texttt{Z3\_VERIFIED} trust
promotion actually buys.

The implementation is split across:
\begin{itemize}[leftmargin=1.7em]
  \item \texttt{deppy/core/z3\_bridge.py} --- low-level wrapper
    around the \texttt{z3} Python package;
  \item \texttt{deppy/core/z3\_encoder.py} --- typed encoding from
    \texttt{SynType} to Z3 sorts;
  \item \texttt{deppy/lean/sidecar\_foundation\_discharge.py} ---
    foundation-level discharge driver;
  \item \texttt{deppy/lean/sidecar\_counterexample\_link.py} ---
    counterexample search driver;
  \item \texttt{deppy/pipeline/counterexample.py} ---
    \texttt{CounterexampleFinder}.
\end{itemize}

\section{When Z3 fires}
\label{sec:z3-when}

The \Deppy{} pipeline calls \Zthree{} in three contexts:

\begin{enumerate}[leftmargin=1.7em]
  \item \textbf{Foundation discharge.}  Each
    \texttt{FoundationDecl} is offered to Z3.  Successful
    discharge promotes the foundation from
    \texttt{AXIOM\_TRUSTED} to \texttt{Z3\_VERIFIED}.

  \item \textbf{PSDL \texttt{assert P, "z3"}.}  The PSDL compiler
    generates a \texttt{Z3Proof(formula=P, binders=\dots)} term;
    \texttt{ProofKernel.verify} dispatches it to the backend.

  \item \textbf{Counterexample search.}  After foundation discharge,
    every formula Z3 \emph{attempted} but did not verify is offered
    to the counterexample finder, which tries to satisfy the
    negation.  When a model is found, the foundation is provably
    false.
\end{enumerate}

A fourth ``passing'' use is when the safety pipeline
(\texttt{deppy/pipeline/safety\_pipeline.py}) discharges per-source
exception-safety obligations; that is conceptually the same
mechanism but lives outside the verify-block path.

\section{Trust promotion}
\label{sec:z3-trust}

The \texttt{TrustLevel} ordering (\Cref{ch:cubical}) places
\texttt{Z3\_VERIFIED} (5) above \texttt{AXIOM\_TRUSTED} (2) and
just below \texttt{KERNEL} (6).  The semantics:

\begin{description}
  \item[\texttt{AXIOM\_TRUSTED}] the proof is admitted on the user's
    authority.  No machine has checked it.
  \item[\texttt{Z3\_VERIFIED}] the SMT solver returned
    \texttt{unsat} on the negation.  We trust Z3's answer.
  \item[\texttt{KERNEL}] every sub-proof was checked by an internal
    rule (or by Lean's kernel via \texttt{LeanProof}).
\end{description}

The promotion from \texttt{AXIOM\_TRUSTED} to
\texttt{Z3\_VERIFIED} is meaningful: the foundation is no longer
admitted by the user; it is checked by an external decision
procedure.  The trust surface shifts from ``the user's claim is
correct'' to ``Z3 is correct on this particular formula''.  We
believe Z3 is significantly more reliable than typical user-supplied
admissions for the formulas it handles.

\section{The typed Z3 encoder}
\label{sec:typed-z3}

\texttt{deppy/core/z3\_encoder.py} translates from the kernel's
\texttt{SynType} universe to Z3 sorts and from \texttt{SynTerm}
expressions to Z3 \texttt{Expr}s.

\subsection{Sorts}

The mapping:
\begin{center}
\begin{tabular}{ll}
\toprule
\textbf{\texttt{SynType}}        & \textbf{Z3 sort} \\
\midrule
\texttt{PyIntType}                & \texttt{Int} \\
\texttt{PyFloatType}              & \texttt{Real} \\
\texttt{PyBoolType}               & \texttt{Bool} \\
\texttt{PyStrType}                & \texttt{String} \\
\texttt{PyListType[T]}            & \texttt{Seq T} (or array of $T$) \\
\texttt{PyDictType[K, V]}         & \texttt{Array K V} \\
\texttt{OptionalType[T]}          & \texttt{Datatype Some T \(\cup\) None} \\
\texttt{PyClassType{n; \dots}}      & uninterpreted sort with arity 0 \\
\texttt{RefinementType{T \mid p}}   & sort of $T$ + assertion of $p$ \\
\texttt{IntervalType[a, b]}       & \texttt{Int} + assertion $a \le x \le b$ \\
\texttt{UnionType[T1, T2]}        & \texttt{Datatype Inl T1 \(\cup\) Inr T2} \\
\bottomrule
\end{tabular}
\end{center}

\subsection{Expressions}

The \texttt{SynTerm} expressions translate directly:
\begin{itemize}[leftmargin=1.7em]
  \item \texttt{Var(name)} → free Z3 constant of the appropriate
    sort.
  \item \texttt{Literal(v)} → Z3 literal.
  \item \texttt{App(f, a)} → \texttt{f(a)} where \texttt{f} is an
    uninterpreted function.
  \item \texttt{Lam} / \texttt{Pi} / path constructors are
    \emph{not} encoded; they are not first-class in Z3.
  \item \texttt{ite(c, t, e)} → \texttt{If(c, t, e)}.
  \item \texttt{Add}/\texttt{Sub}/\texttt{Mul}/\texttt{Div}/\texttt{Mod} → direct.
  \item Comparisons → direct.
\end{itemize}

\subsection{Untyped fallback}

When the kernel cannot infer types (e.g.\ from a free-text formula
without binders), it falls back to a heuristic
\texttt{Int}/\texttt{Bool} encoding: every identifier becomes
\texttt{Int}; comparisons stay; arithmetic stays.  This handles the
canonical arithmetic foundations.  When a formula references
\texttt{sqrt} or \texttt{abs}, the encoder emits an uninterpreted
function symbol; Z3 has no inherent semantics for it, so the
formula is unlikely to discharge.

\section{Foundation discharge in detail}
\label{sec:foundation-discharge-z3}

\texttt{discharge\_foundations(sm)} (file:
\texttt{deppy/lean/sidecar\_foundation\_discharge.py}) is the entry
point.  Per-foundation flow:

\begin{enumerate}[leftmargin=1.7em]
  \item \textbf{Filter.}  \texttt{\_looks\_arithmetic(stmt)} screens
    out statements containing \texttt{sqrt}, \texttt{sin}, \texttt{cos},
    \texttt{Σ}, \texttt{is\_collinear}, \texttt{Triangle},
    \texttt{Point}, etc.  These foundations are pre-classified as
    not arithmetic-Z3-friendly.
  \item \textbf{Strip the guard.}
    \texttt{\_strip\_when\_clause("\dots when \dots")} returns
    \texttt{(core, guard)}.
  \item \textbf{Compose.}
    \texttt{\_rewrite\_for\_z3(stmt)} returns
    \texttt{(guard) => (core)} when there is a guard, else just
    \texttt{core}.
  \item \textbf{Bind.}  \texttt{\_collect\_binders(formula)}
    enumerates free identifiers and assigns each \texttt{int}.
  \item \textbf{Verify.}  Build a \texttt{Z3Proof(formula,
    binders)} and a synthetic \texttt{Judgment}; call
    \texttt{kernel.verify}.
\end{enumerate}

Outcomes for \texttt{sympy\_geometry.deppy}'s 14 foundations:

\begin{center}
\begin{tabular}{ll}
\toprule
\textbf{Foundation} & \textbf{Outcome} \\
\midrule
\texttt{Real\_abs\_nonneg}                       & verified (omega) \\
\texttt{Real\_abs\_zero\_iff\_zero}              & verified (omega) \\
\texttt{Real\_add\_assoc}                         & verified (omega) \\
\texttt{Real\_add\_comm}                          & verified (omega) \\
\texttt{Real\_avg\_minus\_endpoints}              & verified (omega) \\
\texttt{Real\_mul\_div\_self}                     & verified (omega) \\
\texttt{Real\_translate\_killed\_in\_diff}        & verified (omega) \\
\texttt{Real\_scale\_quadratic}                   & admitted (non-linear) \\
\texttt{Real\_sqrt\_nonneg}                       & admitted (sqrt) \\
\texttt{Real\_sqrt\_sq\_nonneg}                   & admitted (sqrt) \\
\texttt{Real\_sqrt\_zero\_iff\_zero}              & admitted (sqrt) \\
\texttt{Real\_sub\_sq\_swap}                      & admitted (non-linear) \\
\texttt{Real\_sub\_sq\_zero\_iff\_eq}             & admitted (non-linear) \\
\texttt{Real\_sum\_of\_nonneg\_zero\_iff\_all\_zero} & admitted (NL prose) \\
\bottomrule
\end{tabular}
\end{center}

The \texttt{omega} tactic in \Lean{}~4 is roughly Z3 over linear
integer arithmetic.  When Z3 succeeds via the encoder, the
generated \Lean{} foundation reads:

\begin{lstlisting}[language=lean4]
theorem Foundation_Real_add_comm : Foundation_Real_add_comm_prop := by
  unfold Foundation_Real_add_comm_prop; intros; omega
\end{lstlisting}

When Z3 cannot handle the formula (non-linear, transcendental, or
prose), the certificate falls back to \texttt{axiom}.

\section{Counterexample search}
\label{sec:counter-search-z3}

\texttt{run\_counterexample\_search(sm,
foundation\_discharge\_report)} (file:
\texttt{deppy/lean/sidecar\_counterexample\_link.py}) runs only on
foundations whose \texttt{z3\_attempted == True} but
\texttt{z3\_verified == False}.

\subsection{The negation game}

For each candidate foundation $\varphi$:
\begin{enumerate}[leftmargin=1.7em]
  \item Build a \texttt{Judgment} whose \texttt{type\_} is the
    \texttt{RefinementType} with \texttt{predicate=}$\varphi$.
  \item Call
    \texttt{CounterexampleFinder.find(goal, timeout\_ms=2000)}.
  \item The finder asks Z3 to satisfy $\neg \varphi$.
  \item If \texttt{sat}, return the model as the counterexample;
    if \texttt{unsat}, conclude $\varphi$ is true (rare ---
    discharge would have already verified it).
  \item If \texttt{unknown}, no counterexample claim made.
\end{enumerate}

\subsection{What the finder returns}

\texttt{CounterexampleResult} carries:
\begin{itemize}[leftmargin=1.7em]
  \item \texttt{found : bool};
  \item \texttt{inputs : dict} --- the satisfying model
    (e.g.\ \verb|{"x": -1, "y": 0}|);
  \item \texttt{explanation : str} --- a human-readable rationale.
\end{itemize}

\subsection{The cascade}

When a foundation is rejected, two things happen in the certificate:
\begin{enumerate}[leftmargin=1.7em]
  \item The foundation's \Lean{} declaration is replaced with a
    rejection notice.
  \item Every theorem citing the rejected foundation is replaced
    with a rejection notice.
  \item Foundations \emph{transitively dependent} on the rejected
    foundation (via \texttt{sm.\_foundation\_deps}) are also
    rejected, by fixed-point closure.
\end{enumerate}

The fixed-point loop is in \texttt{render\_sidecar\_section}:

\begin{lstlisting}[caption={Rejection cascade in
\texttt{render\_sidecar\_section} (paraphrased).}]
rejected_foundations = {o.foundation for o in
                        counterexample_report.outcomes if o.found}
foundation_deps_map = sm._foundation_deps or {}
while True:
    added = False
    for fname, deps in foundation_deps_map.items():
        if fname in rejected_foundations:
            continue
        if any(d in rejected_foundations for d in deps):
            rejected_foundations.add(fname)
            added = True
    if not added:
        break
\end{lstlisting}

\subsection{Why this matters}

A common verifier failure mode is to admit a false foundation and
then claim the theorems it implies are verified.  The
counterexample-search loop closes that hole: any false admission
gets surfaced before the certificate is signed.

The price is performance.  Each candidate foundation gets a
2-second timeout; for a sidecar with a hundred admitted
foundations the loop adds about 200 seconds in the worst case.
Caching mitigates this: repeated runs against unchanged
foundations are essentially free.

\section{The \texttt{Z3Proof} term in \texttt{ProofKernel.verify}}
\label{sec:z3-proofterm}

When the kernel encounters a \texttt{Z3Proof(formula, binders)}
inside a \texttt{ProofTerm} tree, the verification rule fires:

\begin{lstlisting}[caption={\texttt{Z3Proof} verification rule
(paraphrased from the kernel).}]
def verify_z3(self, ctx, goal, term):
    if term._cached_result is not None:
        return term._cached_result
    encoded = self.z3_encoder.encode_formula(
        term.formula, binders=term.binders, ctx=ctx,
    )
    if encoded is None:
        return VerificationResult.fail(
            "Z3 encoder could not handle formula", code="E-Z3-ENC")
    solver = Solver()
    solver.add(Not(encoded))
    result = solver.check()
    if result == unsat:
        verdict = VerificationResult.ok(
            trust=TrustLevel.Z3_VERIFIED,
            message=f"Z3 unsat on negation: {term.formula}")
    elif result == sat:
        verdict = VerificationResult.fail(
            f"Z3 found counterexample: {solver.model()}",
            code="E-Z3-CE")
    else:
        verdict = VerificationResult.fail(
            "Z3 returned unknown", code="E-Z3-UNK")
    term._cached_result = verdict
    return verdict
\end{lstlisting}

The cache field on the \texttt{Z3Proof} dataclass means a
particular formula is asked at most once per session.

\section{Limits of Z3}

We close the chapter by being honest about what Z3 cannot do:

\begin{itemize}[leftmargin=1.7em]
  \item \textbf{Transcendentals.}  \texttt{sqrt}, \texttt{sin},
    \texttt{cos}, \texttt{exp}, \texttt{log} have no Z3 semantics.
    The encoder emits an uninterpreted function and Z3 cannot
    reason about it.

  \item \textbf{Quantifier alternation.}  Z3 accepts $\forall x.\,
    \varphi(x)$ but is incomplete in the presence of nested
    $\forall \exists$.  The arithmetic foundations we discharge are
    universal-only; richer claims are out of scope.

  \item \textbf{Strings and arrays beyond a fixed alphabet.}  Z3's
    string and sequence theories are decidable but \emph{slow}.
    The encoder uses them sparingly and times out conservatively.

  \item \textbf{Inductive reasoning.}  Z3 cannot prove
    $\forall n.\, P(n)$ by induction.  Inductive properties go
    through PSDL's \texttt{induction(\dots)} construct and
    \Lean{}'s elaborator.

  \item \textbf{Higher-order.}  Z3 does not handle $\Pi$/$\Sigma$
    or $\lambda$-terms in the goal.  These go through PSDL's
    quantifier sugar and \Lean.

  \item \textbf{Floating-point semantics.}  \texttt{Real} in Z3 is
    mathematical reals, not IEEE 754.  Floating-point obligations
    are admitted with caveats; the kernel does not currently make
    floating-point claims.
\end{itemize}

The user-facing rule of thumb: \emph{linear integer arithmetic}
discharges; everything else admits.  PSDL's \verb|assert P, "z3"| is
appropriate for arithmetic side conditions in proof bodies; for
deeper claims, use a \verb|by_lean: |\ block or compose with
foundation citations.

\section{Summary}

The \Zthree{} backend:
\begin{itemize}[leftmargin=1.7em]
  \item handles linear integer arithmetic well, transcendentals
    not at all;
  \item promotes successful foundations from
    \texttt{AXIOM\_TRUSTED} to \texttt{Z3\_VERIFIED};
  \item runs counterexample search on the negation of every
    Z3-attempted-but-not-verified foundation;
  \item rejects any theorem citing a counterexample-refuted
    foundation;
  \item caches per-formula results within a run.
\end{itemize}

The next chapter discusses the cohomology layer: what the
intra-procedural cubical analysis, the inter-procedural simplicial
cohomology, and the bicomplex bridge buy --- and what they don't.
